\chapter*{Cara Menggunakan Buku Ini}
\addcontentsline{toc}{chapter}{Cara Menggunakan Buku Ini}
\markboth{Cara Menggunakan Buku Ini}{Cara Menggunakan Buku Ini}

\section*{Anatomi Bab}
\addcontentsline{toc}{section}{Anatomi Bab}

Setiap bab inti memiliki urutan yang sama:
\begin{enumerate}[label=\arabic*.,leftmargin=*]
  \item \textbf{Outcome Pembelajaran dan Sub-CPMK}: kompetensi yang
        harus dicapai dan kode Sub-CPMK yang terkait.
  \item \textbf{Materi Pokok}: peta isi bab.
  \item \textbf{Lima bagian materi}: uraian konsep, definisi, contoh,
        dan ilustrasi.
  \item \textbf{Contoh}: soal berpenyelesaian lengkap.
  \item \textbf{Aktivitas Pembelajaran dan Praktikum}: latihan
        terbimbing dan increment proyek.
  \item \textbf{Latihan dan Refleksi}: soal latihan bertingkat dan
        pertanyaan refleksi.
  \item \textbf{Asesmen}: instrumen dan rubrik yang mengukur Sub-CPMK.
  \item \textbf{Checklist Pencapaian Kompetensi}: bukti yang harus
        ditunjukkan mahasiswa.
  \item \textbf{Rangkuman}: ringkasan dan jembatan ke bab berikutnya.
  \item \textbf{Evaluasi}: soal evaluasi akhir bab dan rubrik singkat.
\end{enumerate}

\section*{Arti Kotak}
\addcontentsline{toc}{section}{Arti Kotak}

\begin{infobox}
Kotak \textbf{Informasi} memuat konsep pendukung, konteks sejarah,
atau keterkaitan dengan bidang lain.
\end{infobox}

\begin{warningbox}
Kotak \textbf{Peringatan} memuat kesalahan umum, jebakan konseptual,
atau batasan yang perlu diperhatikan.
\end{warningbox}

\begin{tipsbox}
Kotak \textbf{Tips} memuat strategi praktis untuk belajar, menulis
kode, atau mengerjakan ujian.
\end{tipsbox}

\begin{reflectionbox}
Kotak \textbf{Refleksi} mengajak mahasiswa menilai pemahaman dan
kebiasaan belajarnya sendiri.
\end{reflectionbox}

\section*{Proyek Mini-Kompilator}
\addcontentsline{toc}{section}{Proyek Mini-Kompilator}

Proyek berkelanjutan pada buku ini adalah mini-kompilator untuk
\textbf{Subset-C}. Spesifikasi bahasa, kode sumber, contoh program,
dan tes tersedia di direktori \file{code/subsetc}. Proyek dibagi
menjadi enam milestone (M0--M5) yang dikerjakan pada bab-bab tertentu.
Setiap bab praktikum membuka dengan tabel artefak: berkas, status awal,
pekerjaan yang harus dilakukan, dan definisi selesai.

Mahasiswa disarankan menyimpan proyek dalam repositori versi dan
melakukan commit pada setiap milestone. Kode yang ditampilkan di buku
adalah kutipan dari berkas nyata; nomor baris pada buku sesuai dengan
nomor baris berkas.

\section*{Dua Jalur Belajar}
\addcontentsline{toc}{section}{Dua Jalur Belajar}

\textbf{Jalur teori} menekankan definisi, pembuktian, dan analisis.
Mahasiswa membaca bab secara berurutan, mengerjakan latihan teori, dan
menggunakan praktikum sebagai ilustrasi. Jalur ini cocok bagi mahasiswa
yang ingin memperdalam dasar formal.

\textbf{Jalur praktik} menekankan implementasi. Mahasiswa membaca bab
secukupnya, langsung mengerjakan praktikum, lalu kembali ke teori ketika
menemui kesulitan. Jalur ini cocok bagi mahasiswa yang belajar paling
baik melalui kode dan eksperimen.

Kedua jalur bertemu pada proyek akhir dan UAS. Tidak ada jalur yang
lebih mudah: keduanya menuntut pemahaman konsep dan kemampuan
implementasi.

\section*{Konvensi Notasi}
\addcontentsline{toc}{section}{Konvensi Notasi}

\begin{itemize}[leftmargin=*]
  \item Simbol nonterminal ditulis miring, misalnya \(E\), \(T\),
        dan \(F\).
  \item Simbol terminal ditulis dengan huruf tebal atau dalam tanda
        kutip, misalnya \code{id}, \code{+}, dan \code{if}.
  \item String kosong dilambangkan dengan \(\varepsilon\).
  \item Turunan ditulis dengan \(\Rightarrow\), dan turunan nol atau
        lebih dengan \(\Rightarrow^{*}\).
  \item Nama berkas dan perintah ditulis dengan huruf mesin, misalnya
        \file{lexer.c} dan \code{pdflatex}.
  \item Istilah Inggris dimiringkan pada kemunculan pertama, misalnya
        \eng{parser}, lalu bentuk Indonesia dipakai konsisten.
\end{itemize}

\section*{Mengompilasi Bab Secara Mandiri}
\addcontentsline{toc}{section}{Mengompilasi Bab Secara Mandiri}

Setiap bab dapat dikompilasi secara mandiri. Dari direktori bab,
jalankan:
\begin{lstlisting}[language=bash]
pdflatex chapter.tex
pdflatex chapter.tex
\end{lstlisting}
Rujukan ke bab lain yang belum ditulis akan tampil sebagai \code{??}.
Ini normal dan tidak mengganggu pemeriksaan tata letak. Untuk
mengompilasi seluruh buku, gunakan \file{scripts/build-all.bat}.
